\subsection{Конечномерное апостериорное информационное состояние}
В задаче управления при неполной наблюдаемости достаточной статистикой истории наблюдений является апостериорное распределение скрытого состояния, а не только его условное среднее. Поэтому в работе отдельно рассматриваются несколько конечномерных представлений апостериорной информации: только апостериорное среднее, среднее с вероятностью скрытого режима и среднее с дополнительной мерой неопределённости. Это позволяет отделить ценность самой фильтрации от ценности отдельных компонентов фильтрованного информационного состояния.

Сравнение проводится при одинаковой линейной форме правила и на одинаковых тестовых траекториях. Тем самым меняется только информационное представление состояния. В выполненных расчётах добавление вероятности режима и простой агрегированной меры неопределённости не даёт дополнительного выигрыша сверх апостериорного среднего. Этот вывод сохраняется и после стандартизации признаков. В нескольких сценариях выигрыш даёт распределительное расширение состояния. Значит, вопрос состоит не просто в том, чтобы заменить наблюдения фильтрованным средним, а в том, какие именно компоненты апостериорной информации следует передавать правилу денежно-кредитной политики.